\chapter{The compiler}\label{ch:compiling}

A program reaches native G17 bytes by one of two routes into a shared backend (\cref{fig:compile-routes}; its stages
are \cref{fig:backend-stages}). Both routes
check the emitted instructions with Apple's decoder at build time, but the program that runs does not use the decoder.

\section{Input routes}\label{ch:source-native-program}

Metal source goes first through Apple's front end, which produces AIR, and the project translates AIR into its own
intermediate representation (IR). The model kernels (the projections, attention, normalization and generation steps)
do not go through Apple's front end; Python builders construct their IR directly. The Metal-source route is tested on a frozen
census of 197 Metal sources that were not written for this compiler. Today's front end and backend compile 128 of them
(MM~\mm{25.196}); the front end refuses 63, reporting the first construct it cannot translate, and Apple's own compiler
rejects 6 (\code{isa/g17-frontend-refusals.json}, field \code{outcomes}). When the census was run on the hardware, 125
compiled (field \code{census\_outcomes}, MM~\mm{25.149}): three sources compile since, two the front end used to refuse
and one the backend did.

% The compiler's two input routes.
\begin{figure}[htbp]
\centering
\begin{tikzpicture}[node distance=5mm and 5mm]
  \node[apple] (msl) {Metal source\\[-1pt]{\footnotesize written outside the project}};
  \node[apple, right=of msl] (fe) {Apple's front end\\[-1pt]{\footnotesize\code{xcrun metal}}};
  \node[apple, right=of fe] (air) {AIR};
  \node[ours, right=of air] (tr) {AIR translation};
  \node[ours, below=10mm of tr] (ir) {The project's IR};
  \node[ours, left=12mm of ir] (py) {Python kernel builders};
  \node[ours, below=of ir, text width=60mm] (be) {Native backend\\[-1pt]{\footnotesize instruction selection, register
    allocation, encoding, tensor lowering}};
  \node[ours, below=of be, text width=60mm] (prog) {Native program\\[-1pt]{\footnotesize G17 bytes and interface description}};
  \node[ours, below left=9mm and -22mm of prog, text width=46mm] (img) {Image author\\[-1pt]{\footnotesize object, metadata, library, archive}};
  \node[nat, below right=9mm and -22mm of prog, text width=46mm] (rt) {Native runtime\\[-1pt]{\footnotesize builds its own resource and launch state}};
  \draw[flow] (msl) -- (fe); \draw[flow] (fe) -- (air); \draw[flow] (air) -- (tr);
  \draw[flow] (tr) -- (ir); \draw[flow] (py) -- (ir);
  \draw[flow] (ir) -- (be); \draw[flow] (be) -- (prog);
  \draw[flow] (prog.south) -- ++(0,-3mm) -| node[pos=0.75, left, tag] {through Metal} (img.north);
  \draw[flow] (prog.south) -- ++(0,-3mm) -| node[pos=0.75, right, tag] {below Metal} (rt.north);
\end{tikzpicture}
\caption{Input routes to a native program. Gray: Apple's; blue: the project's; red: the project's native
runtime.}
\label{fig:compile-routes}
\end{figure}


The backend outputs native instruction bytes together with a description of the
program's buffers, bindings and execution requirements. For the Metal path the image author packages that output into
the object, metadata, library and archive that Metal loads; the below-Metal runtime takes the native bytes and builds
the measured resource and launch state itself (\cref{ch:below-metal-native}).

\section{The running example in the compiler}\label{ch:running-example-through}

The program of \cref{sec:one-program-ir} was compiled and decoded on the CPU, and its bytes are the ones that ran
below Metal (\cref{ch:below-metal-native}).

The IR is SSA over typed values. A function declares its buffers, each with a binding index and an
element type, and a matrix product is a single IR operation, \code{tensor\_matmul}.

\begin{lstlisting}[style=py]
import struct
from agxforge.g17 import cc, ir, scanlink

a = ir.Buffer("A", 1, elem=ir.F16)
b = ir.Buffer("B", 2, elem=ir.F16)
c = ir.Buffer("C", 3, elem=ir.F32)
fn = ir.Function("tensor_example", [a, b, c])
bd = ir.Builder(fn, fn.block("entry"))
bd.tensor_matmul(a, b, c, M=17, N=19, K=16)
i = bd.builtin("threadgroup_position_in_grid")
one = struct.unpack("<I", struct.pack("<f", 1.0))[0]
bd.store_at(c, i, bd.fadd(bd.load(c, i, type=ir.F32), bd.const(one)))
bd.ret()
ir.verify(fn)

program = cc.compile_function(fn)    # 1,232 bytes, 101 instructions; SHA-256 8bb67a3c...c667
image = scanlink.author(program)     # object, metadata, library and archive for the Metal path
\end{lstlisting}

\code{cc.select} maps each IR operation to a recovered G17 form (\cref{tab:running-example-selection}).

\begin{xltabular}{\linewidth}{@{}>{\hsize=0.71\hsize}X>{\hsize=1.29\hsize}X@{}}
\caption{IR operations of the running example and the G17 forms selected for them.}\label{tab:running-example-selection}\\
\toprule
IR operation & G17 form \\
\midrule
\endfirsthead
\tablecontinued{2}
\toprule
IR operation & G17 form \\
\midrule
\endhead
\code{tensor\_matmul} & a tensor body generated by \code{agxforge/g17/tlower.py}
(\cref{ch:tensor-lowering}) \\*
\code{builtin("threadgroup\_position\_in\_grid")} & \op{14059} of the threadgroup position in x \\*
\code{load} & \op{12682} \\*
\code{const} & \op{11842} \\*
\code{fadd} & \op{998} \\*
\code{store\_at} & \op{17229} \\*
\code{ret} & \op{684} \\*
\bottomrule
\end{xltabular}

\code{tlower} turns the 17 × 19 × 16 shape into two tile rows and two tile columns, with the second of
each partial. The lane index, read with \op{14060} into R0L, gives each lane's row \code{m = 4*(l>>4) +
((l>>1)\&3)} and column \code{col = (l\&8) + 4*(l\&1)}, and then \code{m * ld + col} for each leading dimension; this is
the B-row packing of \cref{sec:coordinate-systems}. The addresses sit in R1–R8. The partial-tile masks are built in
16-bit halves (R11H, R12H, R13H) with \op{612} and \op{437}, because the masked forms take a 16-bit mask operand.
The A fragments load into R20–R23 (tile row 0) and R56–R59 (tile row 1, masked), and the B fragments into R68–R71 and
R72–R75 (the masked second column). The fragment loads publish scoreboard slot 0. Since K = 16 is a single slice, there is
one \op{5107} per output tile, and each waits on slot 0. Each tile is stored from its accumulator group (R24, R32, R40
and R48, each eight-aligned) with \op{17257}, or with \op{17258} where the tile is partial.

Each accumulator group starts on a multiple of eight, and each 16-bit fragment on a multiple of four
(\cref{sec:register-classes}). A fragment is read by two MMAs, kept (0) by the first and released (16) by the
second (\cref{sec:release-does}). The epilogue's fp32 add releases both its sources, and the store releases its value
and index, since each is a last use. The tensor body names R0–R8, R11H–R13H and R20–R75; the epilogue's registers, R9
and R16–R18, all fall outside that set. The register count, 76, is the highest register used plus one; the compiler
does not target a count (\cref{sec:capacity-budgets-spills}).

The threadgroup-position read publishes slot 0, the plain load waits on slot 0 for its index and publishes slot 7, and
the fp32 add waits on slot 7. The compiler's checks run on the emitted bytes and reject a wait on a slot that no instruction fills or a read of a late value without its wait
(\cref{ch:synchronization}). \code{agxforge/g17/releasecheck.py} separately rejects a read of a register after its
release (\cref{ch:image-result-validated}).

\code{program.abi()} records what an executor needs. For this program that is ABI version 5; three bindings, A and B
half and read-only and C float and written; an empty constant pool; 101 main instructions; the 24 distinct (opcode,
length) forms used; the special registers read (the lane index and the threadgroup position); and the requirement of
exact launch geometry.

% Inside the backend: the stages between IR and program bytes, and what each one
% refuses or checks.
\begin{figure}[htbp]
\centering
\begin{tikzpicture}[node distance=3.5mm and 9mm,
    st/.style={ours, text width=40mm, align=left, font=\small, minimum height=8mm},
    note/.style={tag, text width=62mm, align=left, anchor=west, font=\scriptsize}]

  \node[st] (ir) {\textbf{IR}\\[-1pt]{\footnotesize SSA over typed values}};
  \node[note, right=of ir] {A function declares its buffers with a binding index and an element
    type; a matrix product is one operation, \code{tensor\_matmul}.};

  \node[st, below=of ir] (sel) {\textbf{\code{cc.select}}\\[-1pt]{\footnotesize instruction selection}};
  \node[note, right=of sel] {Uses a form only if its fields are mapped and its behaviour was
    measured. Anything else is refused by name: "no measured metadata layout for 5 bindings".};

  \node[st, below=of sel] (tl) {\textbf{\code{tlower}}\\[-1pt]{\footnotesize tensor lowering}};
  \node[note, right=of tl] {Generates a product body from shape and types: masks for partial
    tiles, accumulator groups, tensor loads with scoreboard slots and wait masks.};

  \node[st, below=of tl] (alloc) {\textbf{\code{cc.allocate}}\\[-1pt]{\footnotesize register allocation}};
  \node[note, right=of alloc] {R0--R125 with tuple alignment; lifetime modifiers from liveness;
    nothing live around a back edge is released; encoding holes avoided by explicit table.};

  \node[st, below=of alloc] (enc) {\textbf{encoder}};
  \node[note, right=of enc] {Field widths per (opcode, length) form, from the authoring table.};

  \node[st, below=of enc, draw=ink, fill=white] (prog)
    {\textbf{Native program}\\[-1pt]{\footnotesize bytes, \code{contract()}, \code{abi()}}};
  \node[note, right=of prog] {Code identity, entry, bindings and instruction inventory; register
    count, forms, system registers, metadata class and launch.};

  \foreach \a/\b in {ir/sel, sel/tl, tl/alloc, alloc/enc, enc/prog}
    \draw[flow] (\a) -- (\b);

  % The checks every compile runs, to the left.
  \node[apple, text width=30mm, align=left, font=\scriptsize, left=9mm of enc]
    (dec) {Apple's G17 decoder\\[-1pt]{\scriptsize reads back what the encoder wrote}};
  \draw[flow, <->] (enc) -- (dec);
  \node[box, draw=ours, fill=ours!8, text width=30mm, align=left, font=\scriptsize,
        left=9mm of alloc] (chk) {Released-register check\\[-1pt]{\scriptsize and the CPU emulator}};
  \draw[flow, <->] (alloc) -- (chk);
\end{tikzpicture}
\caption{Inside the backend. Each stage refuses what it cannot justify rather than guessing, and every compile
decodes its own output with Apple's decoder.}
\label{fig:backend-stages}
\end{figure}


\section{General selection and allocation rules}\label{ch:selection-allocation-general}

\code{cc.select} uses a form only if its fields are mapped and its behavior
was measured (the form registry, \code{isa/g17-forms.toml} and \code{isa/g17-scalar-isa.toml}). Anything else it refuses with a message
naming what is missing, for example "no measured metadata layout for 5 bindings"; an unmeasured cooperative binding
class is refused the same way (\code{isa/g17-source-execution.json}).

Where the language leaves behavior undefined, the compiler follows Apple's hardware. Its first lowering of
\code{(uint)x} for a negative float returned the magnitude, where Apple's code returns 0. It now lowers the cast to \op{9320}, reproducing Apple's bytes and output (MM~\mm{25.186}, MM~\mm{25.196}; \cref{sec:conversions}).

\code{cc.allocate} assigns physical registers from R0 to R125 under the rules of \cref{ch:register-file}, with tuple
alignment, and sets the lifetime modifiers from liveness (\cref{sec:register-classes,sec:release-does}). A register that is live
around a back edge is never released (\cref{sec:loops-loop-variable-rule}). Encoding holes are avoided through the
explicit tables \code{\_encodable\_dests} and \code{\_encodable\_fma16\_sources} (\cref{sec:encoding-holes}). The
scalar allocator never chooses a register that a tensor body names (\code{\_prepare\_tensor\_composition} in
\code{agxforge/g17/cc.py}).

\section{Tensor lowering}\label{ch:tensor-lowering}

\code{agxforge/g17/tlower.py} generates a matrix-product body from its shape and types (MM~\mm{23.4}, MM~\mm{25.165}).

\begin{enumerate}
\def\labelenumi{\arabic{enumi}.}
\item Read the lane index, and the simdgroup index when there are several. Compute row, column and element addresses
  in scalar code.
\item Build row and column masks for partial tiles. K beyond the matrix is zero-filled by the masked loads, never read
  out of range.
\item Allocate index scalars, fragments and eight-register accumulator groups.
\item Emit tensor loads with scoreboard slots and wait masks. The first K slice uses the from-zero MMA and later slices
  the accumulating one.
\item For a product that accumulates into C, compute from zero and add C afterwards, as Apple's library does.
\item Store the result, with masks on partial tiles.
\end{enumerate}

A product is spread across the GPU by a row grid, several simdgroups, a column grid and split-K, and only the
combinations verified on hardware are admitted (MM~\mm{25.134}, MM~\mm{25.165}). An epilogue in the product's dispatch
can narrow the output to half, requantize it to int8 with per-channel scales or add a residual (MM~\mm{25.165},
MM~\mm{25.183}). Along K, larger shapes run as a runtime loop that carries the accumulators across trips. The loop is
bit-exact from 1 to 255 trips, and at M512 it is 2.25× faster than the unrolled body, which spills (MM~\mm{25.101}),
so MM~\mm{17} rule 30 forbids unrolling the K loop.
